% GENERATED FILE. Edit canonical lessons, then run npm run generate:books.
% canonical-source-hash: fnv1a64:16aadb094b561064
% canonical-lessons: BN-C151-bondho-kora, BN-C151-gora, BN-C151-sarano, BN-C151-jholano, BN-C151-badha

\chapter{To close, To build, To repair, To hang, To tie}
\label{ch:bn-to-close-to-build-to-repair-to-hang-to-tie}
\hlchaptermodality{151}

\begin{chapteropening}
\textbf{By the end of this chapter:} \emph{I can say to close, to build, to repair, to hang and to tie.}

Complete the last lesson of chapter 151: I can say to close, to build, to repair, to hang and to tie.
\end{chapteropening}

\section[\texorpdfstring{\bn{বন্ধ} \bn{করা} (bôndhô kôrā)}{বন্ধ করা (bôndhô kôrā)}]{\bn{বন্ধ} \bn{করা} (bôndhô kôrā) — to close}

\label{lesson:BN-C151-bondho-kora}



\begin{quote}
Before the new one: say the Bengali for to sneeze, then the Bengali for to cough.
\end{quote}

\subsection*{You'll want to know: \bn{বন্ধ} \bn{করা}}
\textbf{\bn{বন্ধ} \bn{করা}} — \emph{bôndhô kôrā} — \textquotedblleft{}to close\textquotedblright{}.

\begin{grammarlens}[title={What you've built}]
One more word for this chapter.
\end{grammarlens}

\subsection*{Your turn}
\begin{itemize}
\raggedright
  \item \emph{Say it:} \emph{bôndhô kôrā}
  \item \emph{Say it:} \emph{bôndhô kôrā}, once more
  \item \emph{Say it:} say \emph{bôndhô kôrā}
  \item \emph{Recall:} say the Bengali for to sneeze, then the Bengali for to cough, then say \emph{bôndhô kôrā} again
\end{itemize}

\begin{tcolorbox}[breakable,colback=teal!4,colframe=teal!35!black,title={Before you move on}]
What does \bn{বন্ধ} \bn{করা} mean? (\textquotedblleft{}To close\textquotedblright{}.) Say it once more.
\end{tcolorbox}
\section[\texorpdfstring{\bn{গড়া} (gôṛā)}{গড়া (gôṛā)}]{\bn{গড়া} (gôṛā) — to build, to shape}

\label{lesson:BN-C151-gora}



\begin{quote}
Before the new one: say the Bengali for to cough, then the Bengali for to close.
\end{quote}

\subsection*{You'll want to know: \bn{গড়া}}
\textbf{\bn{গড়া}} — \emph{gôṛā} — \textquotedblleft{}to build, to shape\textquotedblright{}.

\begin{grammarlens}[title={What you've built}]
One more word for this chapter.
\end{grammarlens}

\subsection*{Your turn}
\begin{itemize}
\raggedright
  \item \emph{Say it:} \emph{gôṛā}
  \item \emph{Say it:} \emph{gôṛā}, once more
  \item \emph{Say it:} say \emph{gôṛā}
  \item \emph{Recall:} say the Bengali for to cough, then the Bengali for to close, then say \emph{gôṛā} again
\end{itemize}

\begin{tcolorbox}[breakable,colback=teal!4,colframe=teal!35!black,title={Before you move on}]
What does \bn{গড়া} mean? (\textquotedblleft{}To build, to shape\textquotedblright{}.) Say it once more.
\end{tcolorbox}
\section[\texorpdfstring{\bn{সারানো} (sārāno)}{সারানো (sārāno)}]{\bn{সারানো} (sārāno) — to repair}

\label{lesson:BN-C151-sarano}



\begin{quote}
Before the new one: say the Bengali for to close, then the Bengali for to build.
\end{quote}

\subsection*{You'll want to know: \bn{সারানো}}
\textbf{\bn{সারানো}} — \emph{sārāno} — \textquotedblleft{}to repair\textquotedblright{}.

\begin{grammarlens}[title={What you've built}]
One more word for this chapter.
\end{grammarlens}

\subsection*{Your turn}
\begin{itemize}
\raggedright
  \item \emph{Say it:} \emph{sārāno}
  \item \emph{Say it:} \emph{sārāno}, once more
  \item \emph{Say it:} say \emph{sārāno}
  \item \emph{Recall:} say the Bengali for to close, then the Bengali for to build, then say \emph{sārāno} again
\end{itemize}

\begin{tcolorbox}[breakable,colback=teal!4,colframe=teal!35!black,title={Before you move on}]
What does \bn{সারানো} mean? (\textquotedblleft{}To repair\textquotedblright{}.) Say it once more.
\end{tcolorbox}
\section[\texorpdfstring{\bn{ঝোলানো} (jholāno)}{ঝোলানো (jholāno)}]{\bn{ঝোলানো} (jholāno) — to hang}

\label{lesson:BN-C151-jholano}



\begin{quote}
Before the new one: say the Bengali for to build, then the Bengali for to repair.
\end{quote}

\subsection*{You'll want to know: \bn{ঝোলানো}}
\textbf{\bn{ঝোলানো}} — \emph{jholāno} — \textquotedblleft{}to hang\textquotedblright{}.

\begin{grammarlens}[title={What you've built}]
One more word for this chapter.
\end{grammarlens}

\subsection*{Your turn}
\begin{itemize}
\raggedright
  \item \emph{Say it:} \emph{jholāno}
  \item \emph{Say it:} \emph{jholāno}, once more
  \item \emph{Say it:} say \emph{jholāno}
  \item \emph{Recall:} say the Bengali for to build, then the Bengali for to repair, then say \emph{jholāno} again
\end{itemize}

\begin{tcolorbox}[breakable,colback=teal!4,colframe=teal!35!black,title={Before you move on}]
What does \bn{ঝোলানো} mean? (\textquotedblleft{}To hang\textquotedblright{}.) Say it once more.
\end{tcolorbox}
\section[\texorpdfstring{\bn{বাঁধা} (bā̃dhā)}{বাঁধা (bā̃dhā)}]{\bn{বাঁধা} (bā̃dhā) — to tie}

\label{lesson:BN-C151-badha}



\begin{quote}
Before the new one: say the Bengali for to repair, then the Bengali for to hang.
\end{quote}

\subsection*{You'll want to know: \bn{বাঁধা}}
\textbf{\bn{বাঁধা}} — \emph{bā̃dhā} — \textquotedblleft{}to tie\textquotedblright{}.

\begin{grammarlens}[title={What you've built}]
One more word for this chapter.
\end{grammarlens}

\subsection*{Your turn}
\begin{itemize}
\raggedright
  \item \emph{Say it:} \emph{bā̃dhā}
  \item \emph{Say it:} \emph{bā̃dhā}, once more
  \item \emph{Say it:} say \emph{bā̃dhā}
  \item \emph{Recall:} say the Bengali for to repair, then the Bengali for to hang, then say \emph{bā̃dhā} again
\end{itemize}

\begin{tcolorbox}[breakable,colback=teal!4,colframe=teal!35!black,title={Before you move on}]
What does \bn{বাঁধা} mean? (\textquotedblleft{}To tie\textquotedblright{}.) Say it once more.
\end{tcolorbox}
